\studySubsection{practice-calc-08}{练习库：高等数学第 8 讲 一元函数积分学的概念与性质}

\problemAnchor{MATH1-CALC-0028}
\begin{problemBox}
\begin{problemMeta}
\textbf{题目编号：} \problemRef{MATH1-CALC-0028} \quad
\textbf{答案：} \answerRef{MATH1-CALC-0028} \quad
\textbf{来源：} Codex 原创 / K079 深度讲义配套练习（非真题）
\end{problemMeta}
\begin{enumerate}
  \item[\textbf{A}] 判断 \(F(x)=\ln(1+x^2)\) 是否为 \(f(x)=2x/(1+x^2)\) 在 \(\mathbb R\) 上的原函数。
  \item[\textbf{B}] 说明 \(f(x)=1/(1+x^2)\) 在 \(\mathbb R\) 上为何必有原函数，并给出一个原函数。
  \item[\textbf{C}] 设 \(f(x)=\lfloor x\rfloor\)。判断它在 \([0,2]\) 上是否 Riemann 可积、是否存在原函数，分别说明理由。
  \item[\textbf{M}] 令 \(F(0)=0\)，且 \(x\ne0\) 时 \(F(x)=x^3\cos(1/x^2)\)。令 \(f=F'\)。证明 \(f\) 在 \(0\) 处不连续，但 \(F\) 确为 \(f\) 在 \(\mathbb R\) 上的原函数。
\end{enumerate}
\end{problemBox}

\problemAnchor{MATH1-CALC-0029}
\begin{problemBox}
\begin{problemMeta}
\textbf{题目编号：} \problemRef{MATH1-CALC-0029} \quad
\textbf{答案：} \answerRef{MATH1-CALC-0029} \quad
\textbf{来源：} Codex 原创 / K080 深度讲义配套练习（非真题）
\end{problemMeta}
\begin{enumerate}
  \item[\textbf{A}] 求 \(\int(e^x+\cos x)\,dx\)。
  \item[\textbf{B}] 求满足 \(F'(x)=\cos x,\ F(\pi/2)=3\) 的函数 \(F\)。
  \item[\textbf{C}] 写出 \(1/[x(x-1)]\) 在 \(\mathbb R\setminus\{0,1\}\) 上的全部原函数，说明三个连通分支上的常数是否必须相同。
  \item[\textbf{M}] 证明：若 \(F,G\) 都是 \(f\) 在区间 \(I\) 上的原函数，则 \(F-G\) 在 \(I\) 上为常数；并说明 \(I\) 不能任意换成非连通集合。
\end{enumerate}
\end{problemBox}

\problemAnchor{MATH1-CALC-0030}
\begin{problemBox}
\begin{problemMeta}
\textbf{题目编号：} \problemRef{MATH1-CALC-0030} \quad
\textbf{答案：} \answerRef{MATH1-CALC-0030} \quad
\textbf{来源：} Codex 原创 / K081 深度讲义配套练习（非真题）
\end{problemMeta}
\begin{enumerate}
  \item[\textbf{A}] 求 \(\int(x^{1/2}+x^{-3})\,dx\)，并写明定义域限制。
  \item[\textbf{B}] 求 \(\int dx/(9+4x^2)\)。
  \item[\textbf{C}] 求 \(\int\left[\frac3{1+9x^2}+\frac1{\sqrt{1-x^2}}\right]dx\)，并写明所讨论的实区间。
  \item[\textbf{M}] 设 \(a\ne0\)，求 \(\int dx/(ax+b)\)，写出绝对值与连通区间条件，并用求导核验。
\end{enumerate}
\end{problemBox}

\problemAnchor{MATH1-CALC-0037}
\begin{problemBox}
\begin{problemMeta}
\textbf{题目编号：} \problemRef{MATH1-CALC-0037} \quad
\textbf{答案：} \answerRef{MATH1-CALC-0037} \quad
\textbf{来源：} Codex 原创 / K088 深度讲义配套练习（非真题）
\end{problemMeta}
\begin{enumerate}
  \item[\textbf{A}] 用 Riemann 和定义求 \(\lim_{n\to\infty}\frac1n\sum_{k=1}^n(1+k/n)\)。
  \item[\textbf{B}] 求 \(\lim_{n\to\infty}\frac2n\sum_{k=1}^n(1+2k/n)^3\)。
  \item[\textbf{C}] 设 \(f\in C[0,1]\)，证明 \(\sum_{k=1}^nf((k/n)^2)(2k-1)/n^2\to\int_0^1f(x)dx\)。
  \item[\textbf{M}] 设 \(f\in C[a,b]\)。写出以每个等分小区间左端点为采样点的 Riemann 和，并证明其极限为 \(\int_a^bf\)。
\end{enumerate}
\end{problemBox}

\problemAnchor{MATH1-CALC-0038}
\begin{problemBox}
\begin{problemMeta}
\textbf{题目编号：} \problemRef{MATH1-CALC-0038} \quad
\textbf{答案：} \answerRef{MATH1-CALC-0038} \quad
\textbf{来源：} Codex 原创 / K089 深度讲义配套练习（非真题）
\end{problemMeta}
\begin{enumerate}
  \item[\textbf{A}] 判断 \(f(x)=\sqrt{1-x^2}\) 在 \([-1,1]\) 上是否 Riemann 可积。
  \item[\textbf{B}] 判断只有有限个跳跃间断点的有界分段函数在闭区间上是否可积，并说明所用条件。
  \item[\textbf{C}] 说明 \(f(x)=1/x\) 在 \([0,1]\) 上不能作为普通 Riemann 积分，并判断对应反常积分是否收敛。
  \item[\textbf{M}] 设 \(T(x)=x\)（\(x\) 有理）、\(0\)（\(x\) 无理）。证明 \(T\) 在 \([0,1]\) 上不可 Riemann 积。
\end{enumerate}
\end{problemBox}

\problemAnchor{MATH1-CALC-0039}
\begin{problemBox}
\begin{problemMeta}
\textbf{题目编号：} \problemRef{MATH1-CALC-0039} \quad
\textbf{答案：} \answerRef{MATH1-CALC-0039} \quad
\textbf{来源：} Codex 原创 / K090 深度讲义配套练习（非真题）
\end{problemMeta}
\begin{enumerate}
  \item[\textbf{A}] 已知 \(\int_0^2(f+g)=4,\ \int_0^2(2f-g)=1\)，求 \(\int_0^2(3f+2g)\)。
  \item[\textbf{B}] 化简 \(\int_{-1}^1f+\int_1^4f+\int_4^2f\)。
  \item[\textbf{C}] 对固定 \(a<b\)，化简 \(\int_a^xf(t)dt+\int_x^bf(t)dt\)，说明 \(x\) 的允许范围。
  \item[\textbf{M}] 计算 \(\int_{-2}^2|x^2-1|dx\)，要求先按全部分界点拆区间。
\end{enumerate}
\end{problemBox}

\problemAnchor{MATH1-CALC-0040}
\begin{problemBox}
\begin{problemMeta}
\textbf{题目编号：} \problemRef{MATH1-CALC-0040} \quad
\textbf{答案：} \answerRef{MATH1-CALC-0040} \quad
\textbf{来源：} Codex 原创 / K091 深度讲义配套练习（非真题）
\end{problemMeta}
\begin{enumerate}
  \item[\textbf{A}] 不计算原函数，估计 \(I=\int_0^2dx/(2+x)\) 的一个常数上下界。
  \item[\textbf{B}] 设 \(f\) 在 \([0,2]\) 上可积，且 \(|f(x)|\le3\)，证明 \(\left|\int_0^2x f(x)dx\right|\le6\)。
  \item[\textbf{C}] 设 \(f\in C[a,b]\)，证明 \(\int_a^b|f(x)|dx=0\) 当且仅当 \(f\equiv0\)。
  \item[\textbf{M}] 设每个 \(f_n\) 都在 \([0,1]\) 上可积，且 \(0\le f_n(x)\le1/n\)，证明 \(\int_0^1f_n(x)dx\to0\)，并指出证明中使用了哪条点态到积分的传递。
\end{enumerate}
\end{problemBox}

\problemAnchor{MATH1-CALC-0041}
\begin{problemBox}
\begin{problemMeta}
\textbf{题目编号：} \problemRef{MATH1-CALC-0041} \quad
\textbf{答案：} \answerRef{MATH1-CALC-0041} \quad
\textbf{来源：} Codex 原创 / K092 深度讲义配套练习（非真题）
\end{problemMeta}
\begin{enumerate}
  \item[\textbf{A}] 求一个 \(\xi\in[0,2]\)，使 \(\int_0^2x\,dx=2\xi\)。
  \item[\textbf{B}] 证明存在 \(\xi\in[0,\pi]\)，使 \(\int_0^\pi f(x)\sin x\,dx=2f(\xi)\)，其中 \(f\in C[0,\pi]\)。
  \item[\textbf{C}] 若 \(f\) 在 \(x_0\) 连续，证明 \(\lim_{h\to0^+}\frac1{2h}\int_{x_0-h}^{x_0+h}f(t)dt=f(x_0)\)。
  \item[\textbf{M}] 设 \(f\in C[0,1]\) 且 \(\int_0^1f=0\)。证明存在 \(\xi\in[0,1]\) 使 \(f(\xi)=0\)，并说明若去掉连续性为何论证会失效。
\end{enumerate}
\end{problemBox}

\problemAnchor{MATH1-CALC-0045}
\begin{problemBox}
\begin{problemMeta}
\textbf{题目编号：} \problemRef{MATH1-CALC-0045} \quad
\textbf{答案：} \answerRef{MATH1-CALC-0045} \quad
\textbf{来源：} Codex 原创 / K096 深度讲义配套练习（非真题）
\end{problemMeta}
\begin{enumerate}
  \item[\textbf{A}] 设 \(F(x)=\int_1^xe^{-t^2}dt\)，求 \(F(1)\) 与 \(F'(x)\)。
  \item[\textbf{B}] 设 \(f(t)=-1\,(|t|<1),1\,(|t|\ge1)\)，求 \(F(x)=\int_0^xf(t)dt\)，并讨论 \(F\) 在 \(\pm1\) 处的连续性与可导性。
  \item[\textbf{C}] 设 \(F(x)=\int_0^x(t^2-1)dt\)，求其单调区间和局部极值。
  \item[\textbf{M}] 若 \(f\) 是连续偶函数，证明 \(F(x)=\int_0^xf(t)dt\) 为奇函数；若 \(f\) 是连续奇函数，证明 \(F\) 为偶函数。
\end{enumerate}
\end{problemBox}

\problemAnchor{MATH1-CALC-0046}
\begin{problemBox}
\begin{problemMeta}
\textbf{题目编号：} \problemRef{MATH1-CALC-0046} \quad
\textbf{答案：} \answerRef{MATH1-CALC-0046} \quad
\textbf{来源：} Codex 原创 / K097 深度讲义配套练习（非真题）
\end{problemMeta}
\begin{enumerate}
  \item[\textbf{A}] 计算 \(\int_0^{\pi/2}\cos x\,dx\)。
  \item[\textbf{B}] 求 \(\frac d{dx}\int_3^x(1+t^4)dt\)。
  \item[\textbf{C}] 求满足 \(y'(x)=1/(1+x^4),y(1)=0\) 的积分形式解，并说明为何无需初等原函数。
  \item[\textbf{M}] 设 \(f\in C[a,b]\)，且 \(G'=f\) 于含 \([a,b]\) 的区间，证明 \(\int_a^bf=G(b)-G(a)\)，要求从变上限积分函数与“原函数相差常数”出发。
\end{enumerate}
\end{problemBox}

\problemAnchor{MATH1-CALC-0047}
\begin{problemBox}
\begin{problemMeta}
\textbf{题目编号：} \problemRef{MATH1-CALC-0047} \quad
\textbf{答案：} \answerRef{MATH1-CALC-0047} \quad
\textbf{来源：} Codex 原创 / K098 深度讲义配套练习（非真题）
\end{problemMeta}
\begin{enumerate}
  \item[\textbf{A}] 求 \(\frac d{dx}\int_1^{\sin x}e^t dt\)。
  \item[\textbf{B}] 求 \(\frac d{dx}\int_{x^2}^{3x}\ln(1+t^2)dt\)。
  \item[\textbf{C}] 设 \(F(x)=\int_{-x}^{2x}t\,dt\)，分别用上下限公式与直接积分求 \(F'(x)\)。
  \item[\textbf{M}] 设 \(F_a(x)=\int_{x-a}^{ax}(1+t)dt\)。求 \(F_a'(0)\)，并求使 \(F_a'(0)=0\) 的全部实参数 \(a\)。
\end{enumerate}
\end{problemBox}

\problemAnchor{MATH1-CALC-0048}
\begin{problemBox}
\begin{problemMeta}
\textbf{题目编号：} \problemRef{MATH1-CALC-0048} \quad
\textbf{答案：} \answerRef{MATH1-CALC-0048} \quad
\textbf{来源：} Codex 原创 / K099 深度讲义配套练习（非真题）
\end{problemMeta}
\begin{enumerate}
  \item[\textbf{A}] 设 \(I(x)=\int_0^1\cos(xt)dt\)，在允许交换求导与积分的条件下求 \(I'(x)\)。
  \item[\textbf{B}] 求 \(I(x)=\int_0^x(x^2+t)dt\) 的导数，要求同时写端点项与积分内偏导项。
  \item[\textbf{C}] 求 \(I(x)=\int_x^{x+1}e^{xt}dt\) 的导数，答案可保留为积分形式。
  \item[\textbf{M}] 写出 Leibniz 扩展公式并说明 \(F,F_x\) 连续等条件在逻辑上承担什么作用；用固定端点积分 \(\int_{-1}^{1}(x+t^2)dt\) 反驳“端点固定所以导数为零”。
\end{enumerate}
\end{problemBox}

\problemAnchor{MATH1-CALC-0049}
\begin{problemBox}
\begin{problemMeta}
\textbf{题目编号：} \problemRef{MATH1-CALC-0049} \quad
\textbf{答案：} \answerRef{MATH1-CALC-0049} \quad
\textbf{来源：} Codex 原创 / K100 深度讲义配套练习（非真题）
\end{problemMeta}
\begin{enumerate}
  \item[\textbf{A}] 求 \(\lim_{n\to\infty}\frac1n\sum_{k=1}^n e^{k/n}\)。
  \item[\textbf{B}] 求 \(\lim_{n\to\infty}\frac1{n^3}\sum_{k=1}^nk(n-k)\)，要求先归一化为 Riemann 和。
  \item[\textbf{C}] 设 \(f\in C[0,2]\)，求 \(\lim_{n\to\infty}\frac2n\sum_{k=1}^nf((2k-1)/n)\)。
  \item[\textbf{M}] 求 \(\lim_{n\to\infty}\frac3n\sum_{k=1}^n\sqrt{1+3k/n}\)，明确写出 \(\Delta x\)、积分区间与采样点。
\end{enumerate}
\end{problemBox}
